\scanpage{5}\mspage{4}

Picard - one can say that the classification ``in the Picard sense'' is finer than the classification in the sense of (Gr)-categories. \deleted{\unsure{x}}

Associating to each $X\handwritten{\in}\mathrm{Ob}\ C$ the symmetry of $X\otimes X$, interpreted as \deleted{\rlap{xxxx}bein} being tensoring by a well-determined element $s(X)$ of $\handwritten{\pi}_{1}(C)=\mathrm{Aut}(1_C)$, one finds a\deleted{e} canonical \deleted{\rlap{xxxxxxxxxxx}map}\inserted{homomorphism}
\[
s(C)=s:\handwritten{\pi}_{\deleted{1}\handwritten{0}}\handwritten{\longrightarrow}\unsure{{}'|\cdot}{}_{2}(\handwritten{\pi}_{\deleted{0}\handwritten{1}})\ ,
\]
where $ {}_{2}(\handwritten{\pi}_{\deleted{0}1})$ denotes the subgroup of elements $x$ of $\handwritten{\pi}_{1}$ such that $2x=0$.

The main result of the structure theory of Picard categories can be stated by saying that $C$ is determined (up to \inserted{\handwritten{noncanonical}} equivalence of Picard categories) by the data of the commutative groups $\handwritten{\pi}_{0}$ and $\handwritten{\pi}_{1}$, and the homomorphism $s:\handwritten{\pi}_{0}\handwritten{\longrightarrow}{}_{2}(\handwritten{\pi}_{1})$. Saying that this homomorphism is zero means moreover that $C$ is a \uline{strict} Picard category, and one recovers the result of \handwritten{[\, ]} according to which a strict Picard category is determined (up to equivalence) by knowledge of its invariants $\handwritten{\pi}_{0}$ and $\handwritten{\pi}_{1}$ alone.

\uline{Remarks}. The theory in the cited work allows us, more precisely, to interpret the \deleted{e}2-category of strict \deleted{E}Picard categories in terms of the full subcategory of the derived category of (Ab) (category of abelian groups), formed by the complexes $K^{\deleted{\unsure{\bullet|\cdot}}}_{\handwritten{\bullet}}$ such that $H^{\deleted{\unsure{i|1}}}_{\handwritten{i}}(K^{\deleted{\bullet}}_{\handwritten{\bullet}})=0$ for $i\handwritten{\ne}0,1$. The preceding result, rather coarse, \inserted{then} reduces to the remark that an object of this category is determined \deleted{\unsure{up to|has}}up to isomorphism (not unique \handwritten{!}) by knowledge of its $H^{\deleted{1}}_{\handwritten{0}}$ $(=\handwritten{\pi}_{0})$ and $H^{\deleted{0}}_{\handwritten{1}}$ $(=\handwritten{\pi}_{1})$. It would be appropriate to develop also for Picard categories not necessarily strict, and for (Gr)-categories, a structure theory for the 2-category they form. Moreover, these \deleted{\rlap{xxxxxxxxx}results}\inserted{developments} still in limbo, \deleted{\unsure{of|this|co}}just like those that \deleted{\rlap{xxxxxxxxxxxxxxxxxxxxxxxx}we have actually} we have just sketched, should one day be \deleted{\rlap{xxxxxxxxxx}extended \unsure{éu|áu|to the}} studied in the framework of stacks (in (Gr)-categories, or in Picard categories) just as in \inserted{\handwritten{Deligne}} \handwritten{[\, ]}.

\deleted{2} \handwritten{3}. \uline{Enveloping Picard categories}.

Let $C$ be an \deleted{\rlap{xx}\unsure{ca|the}} AUC $\otimes$-category (for example a category with a final object and admitting \inserted{\handwritten{\unsure{''|:}}}Cartesian products of two elements), $S$ a subset of $\mathrm{Ob}\ C$. One poses the universal problem of sending (by an AUC $\otimes$-functor) $C$ into a \inserted{$\otimes$-}category \deleted{\unsure{$G$}}\inserted{AUC $G$}, in su\handwritten{\unsure{ch|ch}} a way that for every $X\handwritten{\in} S$, $F(X)$ is an \uline{invertible} object of $G$. The main result here is that this problem always admits a solution, which one may call the AUC $\otimes$-category ``of fractions'' with respect to the set $S$, and \handwritten{\unsure{p|p}}ossibly denote $S^{-1}C$. One finds in fact that for every $G$, the natural functor
% archive sheet number: 4
